% ============================================================================
% ABSTRAK (DRAF): kalimat hasil dan implikasi menunggu E1--E4.
% Batas 150--250 kata, satu paragraf, tanpa rumus, tabel, atau rujukan.
% ============================================================================
\begin{abstract}
Perhitungan PPh 21 harus mengikuti perubahan regulasi dan kebijakan perusahaan, tetapi aturannya umumnya tertanam di kode program atau \en{spreadsheet}. Akibatnya setiap perubahan menuntut pemrograman ulang dan dapat menggeser hasil periode lama tanpa terdeteksi. Penelitian formalisasi hukum pajak berbantuan \en{large language model} (LLM) membangun basis aturan dari nol untuk keadaan hukum pada satu waktu, sehingga belum diketahui apakah pembaruan berulang oleh LLM menjaga hasil periode lama dan apakah kesalahannya merambat dari satu perubahan ke perubahan berikutnya. Penelitian ini mengusulkan \en{knowledge-based system} adaptif yang menyimpan pengetahuan sebagai \en{legal knowledge graph} temporal dua lapis, yaitu lapisan regulasi dan lapisan kebijakan perusahaan, dan menjalankan seluruh perhitungannya dengan \en{inference engine} deterministik. Pembaruan \en{knowledge base} disusun LLM dari dokumen regulasi atau kebijakan, lalu disaring lima gerbang verifikasi dan dipilih lewat konsensus perilaku. Evaluasinya memutar ulang perubahan regulasi PPh 21 pegawai tetap yang nyata pada 2023--2026, termasuk pergantian metode ke tarif efektif rata-rata, menguji perubahan yang belum pernah dilihat perancang, dan membandingkan usulan dengan pembanding serta ablasi. \isian{Kalimat hasil: RA pada perubahan tak terlihat, SER dan CAR beserta selang kepercayaan 95\%, OAR dan kesalahan terpicu pada \en{regulation replay}, serta \en{recall} analisis dampak.} \isian{Kalimat implikasi bagi sistem payroll dan formalisasi peraturan yang sering berubah.}
\end{abstract}

\begin{IEEEkeywords}
Akuisisi pengetahuan, \en{knowledge-based system}, \en{knowledge graph}, \en{large language model}, perubahan regulasi, PPh Pasal 21, \en{rules as code}
\end{IEEEkeywords}

\catatan{Draf ini mengikuti urutan \texttt{paper/draft/skeleton.tex}. Bagian I--III ditulis lengkap. Bagian IV--VI sudah berupa prosa, tetapi memuat keputusan yang belum diambil (kotak oranye) dan butir yang belum didefinisikan (kotak hijau). Abstrak dan Bagian VII--IX hanya kerangka, karena eksperimen belum dijalankan. Kotak ungu menyatakan status yang sebenarnya, yaitu bagian mana yang baru rancangan dan bagian mana yang sudah ada di prototipe. Semua penanda hilang dengan \texttt{\textbackslash rancanganfalse}.}
